% Pursuing Stacks (cotes 134-2 to 134-8): archive cote 134-8, site page 35 (PDF page 36), typescript page 589.
% typescript. Draft compared: ps7.tex:1641-1691.
% AI-transcribed by Claude Opus 5.5 (claude-opus-5-5, adaptive thinking, effort high), Message Batches API, batch ps4 (msgbatch_012Aoc1NuskATPsyBeKSYK7J), 2026-11-12; not yet checked by eye.
\scanpage{134-8/35}
\mspage{589}

they looked just the same as familiar \xout{\unsure{formu}} one's from the standard
duality formulaire (with Lf$_{\hand{!}}$, f$^*$ replaced by Rf$_{\hand{!}}$, f$^{\hand{!}}$)\hand{.} The first
one reads, in case as ab\over{\unsure{xo}}{o}ve when A has a final object \hand{$\varepsilon$}
\[ \text{(2)} \qquad Rf_*(f^*(K)) \overset{\hand{?}}{\simeq} \underline{Hom}_k(f_{\hand{!}}(k_A),K) \quad , \]
\deleted{now in case formation of f$_{\hand{!}}$ commutes \xout{to} to arbitrary base change}
\deleted{(which is true when f is a fibering functor), the right hand side}
\deleted{is isomorphic to \xout{Rx} f$_*$($\underline{\mathrm{Hom}}_k$(k$_A$,f$^*$(K)), and (2) may be rewritten}
\deleted{\[ Rf_*(f^*(K)) \quad f_*(f^*(K)) \quad \text{i.e.} \quad R^if_*(f^*(K))= 0 \text{ for } i \ 0 \]}
\deleted{which should be valid for any sheaf K on B, provided f is fibering and}
\deleted{A has a final object.}
But we have
\[ f_{\hand{!}}(k_A) = f_{\hand{!}}(k^{(\hand{\varepsilon})}) = k^{(b)} \quad , \quad \text{where} \quad b = f_{\hand{!}}\hand{(\varepsilon) = f(\varepsilon)}\hand{,} \]
and if f takes final object into final object, we thus get
f$_{\hand{!}}$(k$_A$) = k$_B$ , and the right hand side of (2) is just K , and hence
(2) implies \deleted{\hand{\illegible{1}}} \xout{R$^i$f$_*$(f$^*$(}
\[ R^if_*(f^*(K)) = 0 \quad \text{for} \quad i \hand{>} 0 , \]
which however needs not be true even for K = k$_B$ , as we saw with the
previous example.

To come back to the general Lf$_{\hand{!}}$, Rf$_*$ formalism, it seems it
can't be helped we'll have to do the usual silly thing and just resol-%
ve the argument involved, projectively in one case, injectively in
the other. Even in the ``absolute case'' when B = e , we could'nt help
either taking such resolutions, when it comes to working with \hand{the}
inter\over{h}{n}al operations \xout{as} R$\underline{\mathrm{Hom}}_k$(L.,K$^\bullet$) or F. $\over{@}{\otimes}^{\hand{L}}_k$ L\unsure{!|.} , as we saw already
in section 108 (which should have tuned down a little my commitedness
to integrators, but it did'nt ! ).

Once this got clear, in order to turn the homology algebra
mill, all we still need is a handy criterion for existence of \inserted[typed]{``enough''} projec-%
tives or injectives in categories of the type
\[ A^{\wedge}_{\hand{\mathcal{M}}} = \underline{Hom}(A^o, \hand{\mathcal{M}}) \quad , \]
with \hand{$\mathcal{M}$} an abelian category. In section 109 we got such a criterion
(prop.4 p.433) - the simplest one can imagine: it is \xout{enough} suffi-%
cient that (besides the stability under direct or inverse limits,
needed anyhow in order for the functor f$_{\hand{!}}$ resp. f$_*$ from A$^{\wedge}_{\hand{\mathcal{M}}}$ to B$^{\wedge}_{\hand{\mathcal{M}}}$

\margin{35}
